\section{Donn\'ees et pr\'eparation du corpus}
\label{sec:donnees}

\subsection{Source et contenu}

Le corpus est constitu\'e de \textbf{36 articles Wikipedia en anglais}, couvrant les principaux
concepts du NLP et du machine learning abord\'es dans le module (mod\`eles de langage, embeddings,
architectures neuronales, m\'etriques d'\'evaluation, etc.). Les articles sont t\'el\'echarg\'es
automatiquement avec la biblioth\`eque \texttt{wikipedia} et sauvegard\'es en texte brut dans
\texttt{data/raw/}.

\subsection{Nettoyage}

Deux articles se sont r\'ev\'el\'es \`etre des pages d'homonymie (\emph{BM25} et
\emph{Tokenization}) sans contenu encyclop\'edique exploitable ; ils ont \'et\'e \'ecart\'es. Le
sujet de la tokenisation est d\'ej\`a couvert par l'article \emph{Tokenization (lexical analysis)}.
Le corpus final compte donc \textbf{34 documents}, ce qui respecte la contrainte minimale de
30 documents impos\'ee par le sujet.

Pour chaque document conserv\'e, les sections de fin g\'en\'eriques (\emph{See also},
\emph{References}, \emph{Notes}, \emph{External links}) sont supprim\'ees afin de ne garder que
le contenu pertinent.

\subsection{D\'ecoupage en chunks}

Chaque document nettoy\'e est d\'ecoup\'e en chunks de \textbf{512 caract\`eres} avec un
chevauchement de \textbf{50 caract\`eres} entre chunks cons\'ecutifs. Ce chevauchement limite le
risque qu'une information importante soit coup\'ee entre deux segments et perdue lors de la
recherche. Le d\'ecoupage produit \textbf{1443 chunks}, stock\'es au format JSON Lines dans
\texttt{data/processed/chunks.jsonl} (identifiant, document source, texte).

\subsection{Liste des documents}

Le tableau~\ref{tab:corpus} liste les 34 articles utilis\'es (titres originaux en anglais).

\begin{table}[H]
  \centering
  \caption{Articles Wikipedia du corpus (34 documents).}
  \label{tab:corpus}
  \small
  \begin{tabular}{@{}p{0.46\textwidth}p{0.46\textwidth}@{}}
    \toprule
    \textbf{Document} & \textbf{Document} \\
    \midrule
    Attention (machine learning) & BERT (language model) \\
    BLEU & Bag-of-words model \\
    Cosine similarity & FastText \\
    Fine-tuning (deep learning) & GPT (language model) \\
    GloVe (machine learning) & Hallucination (artificial intelligence) \\
    Information retrieval & LangChain \\
    Lemmatization & Long short-term memory \\
    N-gram & Named-entity recognition \\
    Perplexity (natural language processing) & Question answering (computing) \\
    ROUGE (metric) & Recurrent neural network \\
    Retrieval-augmented generation & Sentence embedding \\
    Sentiment analysis & Seq2seq \\
    Stemming & Stop word \\
    TF-IDF & Text classification \\
    Text segmentation & Tokenization (lexical analysis) \\
    Transfer learning & Transformer (machine learning) \\
    Vector database & Word2vec \\
    \bottomrule
  \end{tabular}
\end{table}
